% ============================================================
% BME VIK TDK abstract — VARIANT 5: novelty-led (recommended)
%
% Fifth framing, and the one I would submit. It carries the same
% three pieces of work as variant 4, but weights them differently.
%
%   v3 (abstract-v3-novelty.tex)  — novelty alone. One claim, but
%     says "graph-conditioned" without ever earning it: no control
%     separates the graph from the neighbours' text.
%   v4 (abstract-v4-stack.tex)    — all three as equal stages. Airtight
%     logically, but the hook lands in paragraph five and it reads as
%     a roadmap rather than a finding.
%   v5 (this file)                — novelty leads; the calibration and
%     the attribution controls appear as the two conditions that make
%     the residual readable, not as stages of their own.
%
% The structural bet: an abstract is read in its first three sentences,
% and a jury remembers one claim, not three. v5 spends its opening on
% the idea and its middle on why the idea is trustworthy. The controls
% lose top billing but lose none of their content — paragraph four
% still commits to every experiment v4 promised.
%
% Structure follows the official one-page sample:
% https://portal.vik.bme.hu/files/00000047.pdf
%   - max 1 page, min ~20 lines of body text
%   - NO header, footer, or page numbers
%   - figures/equations only if strictly necessary
%   - own contribution must be clearly separated from cited work
%   - body must cover: methodology, key results, conclusions
%
% NUMBERS: two slots. Paragraph four takes the calibration figure
% (share of downstream performance recovered) and the margin over the
% text-matched baseline; paragraph five takes the validation figures.
% Do not submit without them — the official sample requires key
% results in the body.
%
% NOTE ON LANGUAGE: the official TDK conference language is Hungarian
% (English talks require prior agreement with the section chair, per
% the official sample). This English draft is for internal review with
% Bálint/Dani; translate to Hungarian before submission unless an
% English exception is confirmed. Hungarian runs longer — if it
% overflows, cut the third sentence of paragraph two (the
% LLM-as-reviewer line) first; it is the least load-bearing.
% ============================================================

\documentclass[12pt]{article}

\usepackage[utf8]{inputenc}
\usepackage[T1]{fontenc}
\usepackage[english]{babel}
\usepackage[a4paper, margin=2.5cm]{geometry}
\usepackage{mathptmx}      % Times-like font
\usepackage{hyperref}
\usepackage{enumitem}

\pagestyle{empty}          % <- enforces "no header/footer/page numbers"
\setlength{\parindent}{0pt}
\setlength{\parskip}{0.6em}

\begin{document}

% ---------- Title ----------
\begin{center}
    {\Large\bfseries What Its Own References Cannot Anticipate: \\
    A Generative Measure of Scientific Novelty}
\end{center}

\vspace{0.8em}

% ---------- Authors ----------
\begin{center}
    Gema Csanád IV. Inf, gemacsanad@gmail.com \\

\end{center}

\vspace{0.4em}

% ---------- Consultant ----------
\begin{center}
    Consultant: Unyi Dániel
\end{center}

\vspace{1.2em}

% ==========================================================
% BODY
% ==========================================================

A paper's contribution is, by definition, its departure from the work
it cites: anything already contained in its references would not be a
contribution. We turn that observation into a measurement of
scientific novelty. In a citation graph — where each node is a
publication, each edge a citation from one paper to another, and
every node carries that paper's abstract as its text content — we
train a model to predict a paper's abstract from its citation
neighbourhood alone, with the paper's own text withheld entirely, and
take the discrepancy between the predicted and the actual abstract as
an estimate of how much of the paper could not have been anticipated
from its own antecedents.

Novelty is currently measured in two ways, neither conditioned on an
individual paper's own antecedents. Reference-based measures score
how atypical the combination of cited journals is, discarding content
altogether. Language-model measures score the perplexity of a paper's
text, asking whether it is unusual in general rather than unusual
given what this particular paper was built on. A third line asks a
language model to read a paper and judge its novelty, which
reproduces the judge's priors rather than measuring anything. Ours is
conditional and generative: the expectation a paper is scored against
is constructed from that paper's own reference list.

Methodologically, a graph neural network encodes the citation
neighbourhood — the cited papers, and their references in turn,
beyond what could fit in a prompt — into a continuous representation,
which a trained adapter projects into soft prefix tokens injected
into the attention layers of a language model. Generation is
conditioned on this representation together with the plain text of
the directly cited papers. The target's own features are masked and
edges run from cited to citing paper, so nothing published after it
can reach it. The residual is scored separately for research domain,
problem, task and methodology, which a bibliography determines to
very different degrees.

A residual is only as meaningful as the prediction it is taken from,
and two things have to be established before it can be read as
novelty. The first is accuracy where prediction is possible, which we
measure without a language model as judge: generated text is
substituted for the real node attributes, and we ask how much of the
downstream node-classification performance that real text supports
survives the substitution. The second is that the prediction is
structural rather than merely well-informed — against a text-matched
baseline given the same neighbourhood content as plain text within
the same token budget, and against prefixes built from a disabled or
a shuffled neighbourhood, we test whether the graph carries anything
the neighbours' own text does not already convey, which prior
graph-conditioned models, reporting end-task performance only, never
test.

The residual still confounds novelty with incomplete bibliographic
coverage, so it is read relative to matched papers of the same field,
year and observed reference count rather than as an absolute score.
Validation is the substantive difficulty, since novelty has no ground
truth; we approach it without one, through axiomatic properties any
such measure should satisfy, separation of recognised landmark papers
from matched controls, and correlation with the disruption index and
later citation impact computed on the same graph.

Our own contribution is the formulation of novelty as a conditional
generative residual, the graph-conditioned model that produces it,
the two controls that make the residual interpretable, and the
validation. We report where the measure agrees with existing
reference-based and perplexity-based scores and, more informatively,
where it disagrees: the papers whose wording is conventional but
whose content their own antecedents could not have anticipated, and
the reverse.

\end{document}
